\section{Conclusion} \label{sec:con}
In this paper, we present a customized RAG framework and a synthetic fine-tuning dataset of prompt-guided explanations to streamline SVA generation from natural language descriptions. The RAG framework integrates dynamic splitting, HybridRetrieval, and an operator-based rechecking flow for precise context retrieval and assertion validation, while the dataset provides layer-by-layer explanation of each SVA spanning the Boolean, sequence, property, and verification layers. 
We also construct the largest evaluation dataset for NL2SVA to date, with $40$ Verilog designs and $229$ formally verified assertions with detailed annotations.
Evaluation on GPT-4o-mini, CodeX, DeepSeek-V3, and a fine-tuned Qwen2.5-Coder-7B-Instruct model shows that combining our customized RAG framework with prompt-guided fine-tuning delivers marked improvements in syntax correctness and functionality match, enabling robust SVA synthesis with minimal manual effort.

% This paper presents RAG-SVAGen, a retrieval-augmented generation framework designed to improve the quality of SVA generation from natural language property descriptions. To address the challenges of domain-specific syntax and semantics in hardware verification, RAGSVAG integrates three proposed key techniques: a dynamic splitting method, a HybridRetrieval mechanism that combines global semantic retrieval with keyword-guided retrieval, and an SVA operator-based rechecking flow.
% To rigorously evaluate the RAGSVAG framework, we construct the largest evaluation dataset for natural language to SVA translation to date, with 40 Verilog designs and 229 formally verified assertions with detailed annotations. We  develop an automatic evaluation framework that validates syntax correctness and functionality match using FPV tools.
% Experimental results across multiple LLMs show that RAGSVAG substantially improves functionality match baseline LLMs and a related work, advancing the automation of assertion-based hardware verification.
% \newpage
%This paper presents RAGSVAG, a retrieval-augmented generation framework that enhances LLM-based SystemVerilog assertion generation from natural language descriptions. RAGSVAG integrates a dynamic splitting technique for database construction, a HybridRetrieval mechanism combining semantic and keyword-guided retrieval, and an SVA operator-based rechecking flow to refine generated assertions. To support systematic evaluation, we construct the largest known dataset for natural language to SVA translation and develop an automatic verification framework. Experimental results across multiple LLMs show that RAGSVAG substantially improves functionality match baseline LLMs and a related work, advancing the automation of assertion-based hardware verification.



% This paper advances the automation of SVA generation by introducing a comprehensive evaluation dataset, an automatic evaluation framework, and a RAG framework for improving LLM assertion generation. 
% Our evaluation dataset, the largest of its kind to our knowledge, comprises $40$ Verilog designs and $229$ verified assertions, each with detailed annotations, enabling evaluation of LLM-based assertion generation at different levels. 
% The proposed evaluation framework automatically checks the functionality equivalence between the LLM-generated assertions and the reference assertions from our evaluation dataset using the FPV tool. 
% To enhance assertion generation ability of LLM, we construct a RAG database from Verilog textbooks and introduce two text-splitting techniques—static and dynamic—to improve contextual retrieval. The dynamic approach, which separates code and text databases, preserves assertion-related knowledge more effectively, leading to improved LLM performance.
% Despite these advancements, challenges remain in fully automating assertion generation for highly complex or evolving design specifications. Future work can explore improved retrieval mechanisms, adaptive prompting strategies, and deeper integration of formal verification feedback into the generation process. Nonetheless, our contributions—a large-scale, richly annotated dataset, an objective evaluation pipeline, and a RAG-based retrieval framework—provide a strong foundation for advancing LLM-driven assertion generation and its applications in hardware verification.